\section{Markov numbers and degenerations of the projective plane}\label{sec:degensofP2}

The goal of this section is to prove Theorem \ref{thmA:nonplanarlimits}, \ie to show that for any Markov number $d$ there is a smooth limit of degree $d$ curves that is not planar. These limits are constructed as Cartier divisors in degenerations of the plane with isolated singularities. In the process we study some of the basic properties of the log terminal $\QQ$-Gorenstein degenerations of $\PP^2$; \eg we compute their Class groups and show that any ample line bundle is globally generated. To show these divisors are not planar we give an upper bound on their gonality that is less than $d-1$, the gonality of a smooth degree~$d$ plane curve.

We start by considering weighted projective planes that are limits of $\PP^2$. To start we show how the Markov equation arises when considering these weighted projective spaces. Suppose there is a flat family
\[
X \to T
\]
over a pointed curve $0\in T$ such that $X$ is $\QQ$-Gorenstein and for $t\in T$ general $X_t\cong \PP^2$ and $X_0\cong \PP(p,q,r)$ is a weighted projective plane (with $p$, $q$, and $r$ coprime). Then in fact, $(p,q,r)=(a^2,b^2,c^2)$ and $(a,b,c)$ satisfy:
\begin{flalign*}
\text{(the Markov equation)}&&a^2+b^2+c^2=3abc.&&\phantom{\text{(Markov Equation)}}
\end{flalign*}
All solutions to the Markov equation are obtained by successively permuting or performing the \textit{mutation} $(a,b,c) \mto (a,b,3ab - c)$ starting from the minimal solution $(1,1,1)$ \cite{Markov, Mutations}. The first few triples in the Markov tree are
\begin{center}
\scalebox{1.2}{\input{devleming-stapleton_figures/markovtree.tex}}
\end{center}
corresponding to the weighted projective spaces $\PP^2$, $\PP(1,1,4)$, $\PP(1,4,25),\dots$.

The fact that Markov numbers show up when considering $\QQ$-Gorenstein degenerations of $\PP^2$ is a consequence of the constancy of the anticanonical volume $(-K_{X_t})^2$ (Kollár addresses the top self-intersection of a $\QQ$-Cartier divisor $D$ over $T$ in much more generality in \cite[Th.\,11]{KollarConstantVolume}). Setting the anticanonical volume of $\PP(p,q,r)$ equal to $(-K_{\PP^2})^2 = 9$ gives
\[
(p+q+r)^2/pqr = 9.
\]
So $(p+q+r) = 3\sqrt{pqr}$. As $p$, $q$, and $r$ are all coprime, the only possibility is they are all perfect squares. In fact, one can show that for any Markov triple $(a,b,c)$ the singularities of the weighted projective space $\PP(a^2,b^2,c^2)$ can be independently smoothed in a $\QQ$-Gorenstein family which implies that every such weighted projective space is a limit of $\PP^2$ \cite[Cor.\,1.2]{HackProk}. These surfaces were first studied by Manetti in~\cite{Manetti} and are called Manetti surfaces.

\begin{definition}[Manetti surfaces]
Fix a Markov triple $(a,b,c)$. Define
\[
M(a,b,c):= \PP(a^2,b^2,c^2).
\]
Define $M(b,c)$ to be the partial smoothing of the index $a^2$ singularity in $M(a,b,c)$ --- \ie $M(b,c)$ has 2 singularities of index $b^2$ and $c^2$. Likewise define $M(c)$ to be the smoothing of the index $a^2$ and $b^2$ singularities in $M(a,b,c)$.
\end{definition}

\begin{remark}
The singularities appearing on Manetti surfaces are examples of $T$ singularities. In the notation of Hacking and Prokhorov, they are $T_1$ singularities, and the versal $\QQ$-Gorenstein deformation space of such a singularity is one-dimensional (\cf \cite[p.\,4]{HackProk}), so any non-trivial deformation of such a singular point smooths it completely. Throughout this section, we will therefore use the terminology \textit{partial smoothing} of a Manetti surface $M_0$ to mean a one parameter flat family $M$ over a smooth pointed curve $0\in T$ such that $M$ has $\QQ$-Gorenstein singularities, $M_0$ is the fiber over $0 \in T$, and for any singular point in $M_0$ the local deformation of the singularity is either a smoothing or a trivial deformation.
\end{remark}
